\clearpage
\section{Camada formal M4: retirada e devolução compostas}
\label{sec:formal-m4}
M4 liga o estado do frasco (M2.4) ao ciclo de vida do empréstimo (M3) no mesmo
universo formal. A retirada parte de um frasco disponível, sem empréstimo ativo,
e produz um empréstimo \texttt{EM\_USO} com o frasco \texttt{EMPRESTADO}; a
devolução parte de um empréstimo ativo, sempre encerra a custódia e deixa o
frasco em um destino operacionalmente coerente. O status do empréstimo é a
autoridade de custódia e a disponibilidade \texttt{EMPRESTADO} equivale a
exatamente um empréstimo ativo por frasco.

Os predicados de coerência de M2 e a unicidade ativa de M3 são reproduzidos no
modelo integrado e comparados mecanicamente com as origens por
\texttt{tools/formal/withdrawal\_return.mjs}; qualquer drift falha \texttt{alloy-check}.
As guardas de autorização (RBAC, tomador, ciência, TCR), o atraso, o vencimento e
os resultados metrológicos são fatos abstratos das operações.

\begingroup
\small
\raggedright
% Gerado por lcqui-spec-doc; não editar.
\subsection*{Evidência M4 --- retirada e devolução compostas}
M4 reúne o estado do Frasco (M2.4) e o ciclo de vida do Emprestimo (M3) no mesmo universo. A disponibilidade \texttt{EMPRESTADO} equivale a exatamente um empréstimo ativo (\texttt{EM\_USO} ou \texttt{ATRASADO}) por frasco. Retirada e devolução preservam a coerência integrada.
\begin{description}
\item[INV-M4-RETIRADA-APTA-001] RetiradaApta: UNSAT.\newline Escopo: \texttt{check RetiradaApta for 4 but exactly 2 Estado}.
\item[INV-M4-RETIRADA-ATIVO-001] RetiradaCriaAtivo: UNSAT.\newline Escopo: \texttt{check RetiradaCriaAtivo for 4 but exactly 2 Estado}.
\item[INV-M4-RETIRADA-DISPONIBILIDADE-001] RetiradaEmprestado: UNSAT.\newline Escopo: \texttt{check RetiradaEmprestado for 4 but exactly 2 Estado}.
\item[INV-M4-RETIRADA-QUARENTENA-001] QuarentenaBloqueiaRetirada: UNSAT.\newline Escopo: \texttt{check QuarentenaBloqueiaRetirada for 4 but exactly 2 Estado}.
\item[INV-M4-RETIRADA-DESCARTE-001] DescarteBloqueiaRetirada: UNSAT.\newline Escopo: \texttt{check DescarteBloqueiaRetirada for 4 but exactly 2 Estado}.
\item[INV-M4-RETIRADA-VENCIDO-001] VencidoSemAutorizacaoBloqueia: UNSAT.\newline Escopo: \texttt{check VencidoSemAutorizacaoBloqueia for 4 but exactly 2 Estado}.
\item[INV-M4-RETIRADA-ABERTURA-001] AberturaNaRetirada: UNSAT.\newline Escopo: \texttt{check AberturaNaRetirada for 4 but exactly 2 Estado}.
\item[INV-M4-RETIRADA-FRAME-001] RetiradaNaoInterfere: UNSAT.\newline Escopo: \texttt{check RetiradaNaoInterfere for 4 but exactly 2 Estado}.
\item[FRAME-M4-RETIRADA-STATUS-001] RetiradaNaoInterfereStatus: UNSAT.\newline Escopo: \texttt{check RetiradaNaoInterfereStatus for 4 but exactly 2 Estado}.
\item[INV-M4-DEVOLUCAO-ATIVO-001] DevolucaoExigeAtivo: UNSAT.\newline Escopo: \texttt{check DevolucaoExigeAtivo for 4 but exactly 2 Estado}.
\item[INV-M4-DEVOLUCAO-CUSTODIA-001] DevolucaoEncerraCustodia: UNSAT.\newline Escopo: \texttt{check DevolucaoEncerraCustodia for 4 but exactly 2 Estado}.
\item[INV-M4-DEVOLUCAO-NORMAL-001] DevolucaoNormal: UNSAT.\newline Escopo: \texttt{check DevolucaoNormal for 4 but exactly 2 Estado}.
\item[INV-M4-DEVOLUCAO-ATRASO-001] DevolucaoAtraso: UNSAT.\newline Escopo: \texttt{check DevolucaoAtraso for 4 but exactly 2 Estado}.
\item[INV-M4-DEVOLUCAO-ANOMALIA-001] DevolucaoAnomalia: UNSAT.\newline Escopo: \texttt{check DevolucaoAnomalia for 4 but exactly 2 Estado}.
\item[INV-M4-ANOMALIA-PRECEDENCIA-001] AnomaliaPrecedeAtraso: UNSAT.\newline Escopo: \texttt{check AnomaliaPrecedeAtraso for 4 but exactly 2 Estado}.
\item[INV-M4-ANOMALIA-RETENCAO-001] AnomaliaRetem: UNSAT.\newline Escopo: \texttt{check AnomaliaRetem for 4 but exactly 2 Estado}.
\item[INV-M4-ESGOTAMENTO-001] EsgotamentoVazio: UNSAT.\newline Escopo: \texttt{check EsgotamentoVazio for 4 but exactly 2 Estado}.
\item[INV-M4-DESTINO-QUARENTENA-001] DestinoQuarentena: UNSAT.\newline Escopo: \texttt{check DestinoQuarentena for 4 but exactly 2 Estado}.
\item[INV-M4-DESTINO-DISPONIVEL-001] DestinoDisponivel: UNSAT.\newline Escopo: \texttt{check DestinoDisponivel for 4 but exactly 2 Estado}.
\item[INV-M4-DESTINO-DESCARTE-001] DestinoDescarte: UNSAT.\newline Escopo: \texttt{check DestinoDescarte for 4 but exactly 2 Estado}.
\item[INV-M4-DEVOLUCAO-FRAME-001] DevolucaoNaoInterfere: UNSAT.\newline Escopo: \texttt{check DevolucaoNaoInterfere for 4 but exactly 2 Estado}.
\item[FRAME-M4-DEVOLUCAO-STATUS-001] DevolucaoNaoInterfereStatus: UNSAT.\newline Escopo: \texttt{check DevolucaoNaoInterfereStatus for 4 but exactly 2 Estado}.
\item[INV-M4-COERENCIA-001] OperacoesPreservamCoerencia: UNSAT.\newline Escopo: \texttt{check OperacoesPreservamCoerencia for 4 but exactly 2 Estado}.
\item[WIT-M4-RETIRADA-ABERTO-001] RetiradaAberto: SAT.\newline Escopo: \texttt{run RetiradaAberto for 4 but exactly 2 Estado}.
\item[WIT-M4-RETIRADA-FECHADO-001] RetiradaFechado: SAT.\newline Escopo: \texttt{run RetiradaFechado for 4 but exactly 2 Estado}.
\item[WIT-M4-RETIRADA-ABERTURA-001] RetiradaAbertura: SAT.\newline Escopo: \texttt{run RetiradaAbertura for 4 but exactly 2 Estado}.
\item[WIT-M4-RETIRADA-EXCEPCIONAL-001] RetiradaExcepcional: SAT.\newline Escopo: \texttt{run RetiradaExcepcional for 4 but exactly 2 Estado}.
\item[WIT-M4-DEVOLUCAO-NORMAL-001] DevolucaoNormalHabitavel: SAT.\newline Escopo: \texttt{run DevolucaoNormalHabitavel for 4 but exactly 2 Estado}.
\item[WIT-M4-DEVOLUCAO-ATRASADA-001] DevolucaoAtrasadaHabitavel: SAT.\newline Escopo: \texttt{run DevolucaoAtrasadaHabitavel for 4 but exactly 2 Estado}.
\item[WIT-M4-DEVOLUCAO-ANOMALA-001] DevolucaoAnomalaHabitavel: SAT.\newline Escopo: \texttt{run DevolucaoAnomalaHabitavel for 4 but exactly 2 Estado}.
\item[WIT-M4-DEVOLUCAO-VAZIO-001] DevolucaoVazioHabitavel: SAT.\newline Escopo: \texttt{run DevolucaoVazioHabitavel for 4 but exactly 2 Estado}.
\item[WIT-M4-DEVOLUCAO-ANOMALIA-ATRASO-001] DevolucaoAnomaliaAtraso: SAT.\newline Escopo: \texttt{run DevolucaoAnomaliaAtraso for 4 but exactly 2 Estado}.
\item[WIT-M4-DEVOLUCAO-ANOMALIA-VAZIO-001] DevolucaoAnomaliaVazio: SAT.\newline Escopo: \texttt{run DevolucaoAnomaliaVazio for 4 but exactly 2 Estado}.
\item[WIT-M4-DEVOLUCAO-VENCIDO-QUARENTENA-001] DevolucaoVencidoQuarentena: SAT.\newline Escopo: \texttt{run DevolucaoVencidoQuarentena for 4 but exactly 2 Estado}.
\item[WIT-M4-DEVOLUCAO-VENCIDO-DISPONIVEL-001] DevolucaoVencidoDisponivel: SAT.\newline Escopo: \texttt{run DevolucaoVencidoDisponivel for 4 but exactly 2 Estado}.
\item[WIT-M4-DEVOLUCAO-VENCIDO-DESCARTE-001] DevolucaoVencidoDescarte: SAT.\newline Escopo: \texttt{run DevolucaoVencidoDescarte for 4 but exactly 2 Estado}.
\item[WIT-M4-DOIS-FRASCOS-001] DoisFrascosAtivos: SAT.\newline Escopo: \texttt{run DoisFrascosAtivos for 4 but exactly 2 Estado}.
\item[WIT-M4-CICLO-001] CicloCompleto: SAT.\newline Escopo: \texttt{run CicloCompleto for 5 but exactly 3 Estado}.
\item[INV-M4-RETIRADA-ATIVO-006] RetiradaCriaAtivoAmpliado: UNSAT.\newline Escopo: \texttt{check RetiradaCriaAtivoAmpliado for 6 but exactly 2 Estado}.
\item[INV-M4-RETIRADA-QUARENTENA-006] QuarentenaBloqueiaRetiradaAmpliado: UNSAT.\newline Escopo: \texttt{check QuarentenaBloqueiaRetiradaAmpliado for 6 but exactly 2 Estado}.
\item[INV-M4-DEVOLUCAO-FRAME-006] DevolucaoNaoInterfereAmpliado: UNSAT.\newline Escopo: \texttt{check DevolucaoNaoInterfereAmpliado for 6 but exactly 2 Estado}.
\item[INV-M4-COERENCIA-006] OperacoesPreservamCoerenciaAmpliado: UNSAT.\newline Escopo: \texttt{check OperacoesPreservamCoerenciaAmpliado for 6 but exactly 2 Estado}.
\item[WIT-M4-CICLO-006] CicloCompletoAmpliado: SAT.\newline Escopo: \texttt{run CicloCompletoAmpliado for 6 but exactly 3 Estado}.
\end{description}
UNSAT para check: nenhum contraexemplo encontrado no escopo declarado. SAT para run: testemunha encontrada. As buscas limitadas não são prova universal.

Não certifica Q06 ou tara (M6), extravio/reencontro/quarentena completos (M5), idempotência transacional (M7), RBAC (M9), Firestore/backend nem concorrência real. A validade calculada, a primeira abertura que vence e os destinos de vencido são representados por fatos abstratos.

\par
\endgroup

\subsection*{Regras compostas verificadas}
A retirada exige disponibilidade, estado físico utilizável, ausência de
quarentena, ausência de pendência/autorização de descarte e ausência de
empréstimo ativo; frasco vencido só é retirado com uso vencido autorizado, o que
mantém \emph{pendente de descarte} bloqueado mesmo com disponibilidade
\texttt{DISPONIVEL}. A primeira abertura durante a retirada leva \texttt{FECHADO}
a \texttt{ABERTO}. A devolução sempre encerra a custódia: sem anomalia e sem
atraso resulta \texttt{DEVOLVIDO}, com atraso \texttt{DEVOLVIDO\_COM\_ATRASO} e
com anomalia \texttt{DEVOLVIDO\_COM\_ANOMALIA}; a anomalia tem precedência sobre
o atraso e retém o frasco em quarentena indisponível com consumo pendente. O
esgotamento confirmado leva a \texttt{VAZIO} indisponível, com precedência sobre
a disponibilidade autorizada; os destinos de vencido (\texttt{QUARENTENA},
\texttt{DISPONIVEL\_AUTORIZADO}, \texttt{PENDENTE\_DE\_DESCARTE}) são coerentes e
o pendente de descarte não pode ser retirado.

\subsection*{Limites e proveniência}
Os fragmentos são vinculados ao IR e às evidências M0--M3 pelo
\texttt{generated/MANIFEST.json}. M4 NÃO calcula Q06, tara, densidade ou
validade (M6); NÃO modela extravio/reencontro nem quarentena manual completos
(M5); NÃO formaliza idempotência transacional (M7), RBAC (M9), Firestore nem
concorrência real. A primeira abertura que vence é tratada como abstração de
validade: o caminho documental exige uso vencido autorizado prévio ou abertura
prévia com destino, não havendo deadlock. Não antecipa M5.
\par
